\documentclass[12pt]{article}
\usepackage[a4paper,margin=1in]{geometry}\usepackage[T1]{fontenc}
\usepackage{lmodern,microtype}\usepackage{amsmath,amssymb,amsthm,mathtools}
\usepackage{booktabs,array,enumitem,xcolor}\usepackage[numbers,sort&compress]{natbib}
\usepackage[colorlinks=true,linkcolor=blue!55!black,citecolor=blue!55!black,urlcolor=blue!55!black]{hyperref}
\linespread{1.07}
\newtheorem{theorem}{Theorem}[section]\newtheorem{proposition}[theorem]{Proposition}
\theoremstyle{remark}\newtheorem{remark}[theorem]{Remark}
\newcommand{\norm}[1]{\left\lVert#1\right\rVert}

\title{Endpoint-Determined Procrustes Maps for Spectral Shells\\
\large Exact Existence, Optimal Cost, and the Predictivity Gap}
\author{Liang Wang}\date{July 2026}
\begin{document}\maketitle
\begin{abstract}
A poor Haar transport angle does not prevent an exact map between two
equal-dimensional endpoint packets: one can always construct a partial
isometry after seeing both spaces.  This paper makes that statement precise
and quantifies its cost for the physical edge quartets.

For orthonormal packet frames $Q_c,Q_f$ and an ambient embedding $J$, the
unitary polar factor of $Q_f^*JQ_c$ solves the orthogonal Procrustes problem.
The resulting rank-four operator $H=Q_fUQ_c^*$ maps the coarse packet exactly
onto the fine packet, satisfies $H^*H=Q_cQ_c^*$ and $HH^*=Q_fQ_f^*$, and
minimizes the discrepancy from the embedded coarse frame.  Its squared
minimum cost is $2k-2\sum_j\cos\theta_j$.

An 80-case identity audit has zero failures and maximum error
$3.15\times10^{-15}$.  For the four physical transitions the rank-normalized
optimal costs range from $0.44075$ to $0.70099$.  Moreover, exact range
transport does not intertwine dynamics: the largest matched eigenvalue
movement is $0.38160$.  The construction is therefore a canonical
endpoint alignment but not a predictive renormalization map, because it
uses the fine packet it is meant to predict.
\end{abstract}

\section{Existence is not yet transport theory}

RH-204 provides a unique correspondence between the two conjugate branches
at adjacent scales \citep{WangRH204}.  It is tempting to declare that this
already gives a shell map.  Linear algebra shows that an exact map indeed
exists, but also shows why that fact alone is weak: any two subspaces of the
same finite dimension admit a partial isometry.

The useful questions are instead:
\begin{enumerate}[leftmargin=2em]
\item which exact map is closest to the predeclared physical embedding;
\item how large is that correction;
\item does the map intertwine the compressed dynamics;
\item can it be constructed from the coarse level alone?
\end{enumerate}

\section{Packet frames and principal cosines}

Let $Q_c\in\mathbb C^{n_c\times k}$ and
$Q_f\in\mathbb C^{n_f\times k}$ have orthonormal columns, and let
$J:\mathbb C^{n_c}\to\mathbb C^{n_f}$ be an isometry.  Set
\begin{equation}\label{eq:overlap}
 M=Q_f^*JQ_c.
\end{equation}
The singular values $s_j$ of $M$ are the principal cosines between
$\operatorname{Ran}(JQ_c)$ and $\operatorname{Ran}(Q_f)$
\citep{StewartSun1990}.  Write
\begin{equation}\label{eq:svd}
 M=L\Sigma R^*,\qquad U=LR^*.
\end{equation}

\section{Optimal endpoint map}

\begin{theorem}[Procrustes shell map]\label{thm:procrustes}
The unitary $U$ in \eqref{eq:svd} minimizes
\begin{equation}\label{eq:minimization}
 \min_{Z^*Z=I_k}\norm{Q_fZ-JQ_c}_F.
\end{equation}
The minimum obeys
\begin{equation}\label{eq:cost}
 \min\norm{Q_fZ-JQ_c}_F^2
 =2k-2\sum_{j=1}^k s_j.
\end{equation}
Moreover,
\begin{equation}\label{eq:H}
 H=Q_fUQ_c^*
\end{equation}
is a partial isometry with
\begin{equation}\label{eq:projectors}
 H^*H=Q_cQ_c^*,\qquad HH^*=Q_fQ_f^*.
\end{equation}
\end{theorem}

\begin{proof}
Expanding the squared norm gives
\begin{equation*}
 2k-2\operatorname{Re}\operatorname{tr}(Z^*M).
\end{equation*}
Von Neumann's trace inequality bounds the last real part by
$\sum_js_j$, with equality for $Z=LR^*$.  Substituting \eqref{eq:H} and
using orthonormality gives \eqref{eq:projectors}.
\end{proof}

We normalize the reported cost by $\sqrt{k}$, so it measures the root mean
square frame correction per mode.

\section{Physical Procrustes costs}

For the four adjacent physical cases, the optimal costs are:
\begin{center}
\begin{tabular}{llrrr}
\toprule
step & side & right cost & left cost & max $|\Delta\lambda|$\\
\midrule
$0.04\to0.02$ & left  & $0.48343$ & $0.45464$ & $0.38105$\\
$0.04\to0.02$ & right & $0.45582$ & $0.51047$ & $0.38160$\\
$0.02\to0.01$ & left  & $0.69446$ & $0.44075$ & $0.13643$\\
$0.02\to0.01$ & right & $0.44486$ & $0.70099$ & $0.13458$\\
\bottomrule
\end{tabular}
\end{center}
The endpoint correction is never small, and refinement does not improve
both the right and left costs simultaneously.  The asymmetry exchanges
between channels on the fine step, a signature of nonnormal dual geometry.

\section{Exact range map versus dynamical intertwining}

Suppose the branch correspondence orders eigenvector frames $V_c,V_f$ and
diagonal eigenvalue matrices $\Lambda_c,\Lambda_f$.  Any map $T$ satisfying
$TV_c=V_f$ obeys
\begin{equation}\label{eq:dynamics}
 (A_fT-TA_c)V_c=V_f(\Lambda_f-\Lambda_c).
\end{equation}
Thus exact range transport cannot make the dynamical defect vanish unless
the matched eigenvalues agree.  In a conditioned norm, the right side gives
an unavoidable spectral movement term.

The largest movement is $0.38160$ on the first transition.  It falls to
about $0.136$ on the second, but remains a visible floor.  The Procrustes map
optimizes geometry, not dynamics.

\section{Right and left maps are distinct}

For a normal operator one might use one unitary map for both right and left
spaces.  Here the packet projectors are oblique.  Applying
Theorem~\ref{thm:procrustes} separately gives $H_R$ and $H_L$, with different
costs.  Nothing in the theorem ensures
\begin{equation*}
 H_L^*H_R=I
\end{equation*}
on biorthogonal packet coordinates.  Enforcing that condition would be an
oblique, not orthogonal, Procrustes problem and would inherit the large
cross-Gram conditioning measured in RH-197.

\section{An eigenvector-interpolating alternative}

Once matched biorthogonal eigenvector frames $V_c,W_c$ and $V_f,W_f$ are
known, one may define on the coarse packet
\begin{equation}\label{eq:interpolating-map}
 T_R=V_fW_c^*,\qquad T_L=W_fV_c^*.
\end{equation}
These maps carry individual labeled right and left vectors exactly when the
frames are biorthonormal.  They make the spectral defect in
\eqref{eq:dynamics} transparent, but their norms inherit oblique
conditioning and they are even more endpoint dependent than the orthogonal
Procrustes map.

Thus there are at least three distinct optimization targets: smallest
orthogonal frame correction, exact labeled-eigenvector interpolation, and
smallest dynamical intertwining defect.  They coincide only in special
normal, spectrally stationary settings.  The physical packets do not satisfy
those hypotheses.

\section{Identity audit}

Eighty random complex tests use packet ranks two through five, rectangular
ambient isometries, and independent endpoint frames.  We check both
identities in \eqref{eq:projectors} and equality of the computed residual
with \eqref{eq:cost}.  There are zero failures at tolerance $10^{-10}$; the
maximum error over all recorded identities is $3.1415\times10^{-15}$.

\section{The predictivity gap}

The map $H$ is endpoint determined: computing $Q_f^*JQ_c$ requires knowing
the fine spectral packet.  Therefore it cannot prove that the fine packet
exists, predict where it lies, or define an inductive limit without prior
spectral information.

This distinguishes two notions:
\begin{description}[leftmargin=2.5em,style=nextline]
\item[Retrospective transport] align two already computed endpoint packets;
\item[Predictive renormalization] construct the fine packet and its error
from coarse data and uniform analytic bounds.
\end{description}
RH-205 proves the former and explicitly leaves the latter open.

\section{Four levels of canonicity}

It is useful to reserve separate language for:
\begin{enumerate}[leftmargin=2em]
\item existence of some finite packet isomorphism;
\item the optimal endpoint map relative to a declared ambient embedding;
\item a map determined from coarse data and the scale parameter;
\item a uniform map producing a convergent inductive system.
\end{enumerate}
Theorem~\ref{thm:procrustes} reaches level two.  Gate A requires substantially
more than level three because the packet rank must eventually grow.

\section{Next invariant: physical residues}

The branch labels permit one to compare residues without choosing an
eigenvector phase.  RH-206 asks whether a single scale-dependent scalar can
renormalize all four residues or whether a genuinely branch-dependent
cocycle is forced.  This test is independent of the endpoint partial
isometry.

\section{Claim boundary}

The Procrustes theorem and partial-isometry identities are exact.  The four
physical costs are finite floating measurements.  No intrinsic coarse-only
map, uniform scale estimate, inductive-limit Hilbert space, Gate-A
determinant, or Hilbert--P\'olya consequence is obtained.

\bibliographystyle{plainnat}\bibliography{references}\end{document}
